% Part: many-valued-logic
% Chapter: syntax-and-semantics
% Section: connectives

\documentclass[../../../include/open-logic-section]{subfiles}

\begin{document}

\olfileid{mvl}{syn}{con}

\olsection{ভাষা ও সংযোজক}

ধ্রুপদি বচনমূলক যুক্তিবিদ্যা এবং আরও অনেক যুক্তিবিদ্যায়
\emph{বচনধ্রুবক} ও \emph{সংযোজক}-এর একটি নির্দিষ্ট ভাণ্ডার ব্যবহার করা
হয়। যেমন, আমরা নিচেরগুলিকে মৌলিক সংকেত হিসেবে ব্যবহার করি:
\begin{enumerate}
\tagitem{prvFalse}{!!{falsity}-র বচনধ্রুবক~$\lfalse$।}{}
\tagitem{prvTrue}{!!{truth}-এর বচনধ্রুবক~$\ltrue$।}{}
\item যৌক্তিক সংযোজকগুলি:
  \startycommalist
  \iftag{prvNot}{\ycomma $\lnot$ (নঞর্থকরণ)}{}%
  \iftag{prvAnd}{\ycomma $\land$ (সংযোজন)}{}%
  \iftag{prvOr}{\ycomma $\lor$ (বিয়োজন)}{}%
  \iftag{prvIf}{\ycomma $\lif$ (!!{conditional})}{}%
  \iftag{prvIff}{\ycomma $\liff$ (!!{biconditional})}{}%
\end{enumerate}
\iftag{defNot,defOr,defAnd,defIf,defIff,defTrue,defFalse,defEx,defAll}{%
উপরের মৌলিক সংযোজকগুলির পাশাপাশি আমরা সংক্ষিপ্ত রূপ হিসেবে সংজ্ঞায়িত
সংকেতও ব্যবহার করি, যেমন
\startycommalist
  \iftag{defNot}{\ycomma $\lnot$ (নঞর্থকরণ)}{}%
  \iftag{defAnd}{\ycomma $\land$ (সংযোজন)}{}%
  \iftag{defOr}{\ycomma $\lor$ (বিয়োজন)}{}%
  \iftag{defIf}{\ycomma $\lif$ (!!{conditional})}{}%
  \iftag{defIff}{\ycomma $\liff$ (!!{biconditional})}{}%
  \iftag{defFalse}{\ycomma $\lfalse$ (!!{falsity})}{}%
  \iftag{defTrue}{\ycomma $\ltrue$ (!!{truth})}.}{}

বহুমানী যুক্তিবিদ্যাতেও একই সংযোজকগুলি ব্যবহৃত হয়। কিন্তু একই
যুক্তিবিদ্যায়, যেমন, সংযোজনের একাধিক সংস্করণ রাখা প্রায়ই উপযোগী; তখন
সংস্করণগুলিকে আলাদা রাখতে পৃথক সংকেত দরকার। কোনো কোনো বহুমানী
যুক্তিবিদ্যায় এমন সংযোজকও থাকে, যার ধ্রুপদি যুক্তিবিদ্যায় কোনো সমতুল্য
নেই। তাই আমরা প্রচলিত রীতির চেয়ে একটু বেশি সাধারণ হব।

\begin{defn}
  একটি \emph{বচনমূলক ভাষা} সংযোজকসমূহের একটি সেট~$\Lang L$ নিয়ে গঠিত।
  প্রতিটি সংযোজক~$\star$-এর একটি \emph{স্থানসংখ্যা} থাকে; স্থানসংখ্যা~$n$
  হলে সংযোজকটিকে \emph{$n$-স্থানীয়} বলা হয়। স্থানসংখ্যা~$0$-এর
  সংযোজকগুলিকে \emph{ধ্রুবক}-ও বলা হয়; স্থানসংখ্যা~$1$-এর সংযোজক
  \emph{এক-স্থানীয়}, আর স্থানসংখ্যা~$2$-এর সংযোজক \emph{দ্বি-স্থানীয়}।
\end{defn}

\begin{ex}
  বচনমূলক যুক্তিবিদ্যার মানক ভাষা~$\Lang L_0$ নিচের সংযোজকগুলি নিয়ে
  গঠিত (সঙ্গে তাদের স্থানসংখ্যা দেওয়া হয়েছে):
  $\lfalse$~($0$),
  $\lnot$~($1$),
  $\land$~($2$),
  $\lor$~($2$),
  $\lif$~($2$)। আমরা যে অধিকাংশ যুক্তিবিদ্যা বিবেচনা করব, সেগুলিতে এই
  ভাষাই ব্যবহৃত হবে। প্রথা ও রীতি অনুসারে কোনো কোনো যুক্তিবিদ্যায় কিছু
  সংযোজকের জন্য আলাদা সংকেত ব্যবহৃত হয়। যেমন, গুণন যুক্তিবিদ্যায়
  সংযোজনের সংকেত প্রায়ই~$\land$-এর বদলে~$\odot$ হয়। কখনো নতুন অপারেটর
  যোগ করা সুবিধাজনক; যেমন নির্ধারিততা অপারেটর~$\triangle$
  ($1$-স্থানীয়)।
\end{ex}

\end{document}
